% Generated by scripts/assemble.py; corrections belong in source/corrections.json.
\examhead{109}{合肥工业大学 2026 年数学分析试卷}{107}
% p108-b08; source PDF page 108
\par\Needspace{5\baselineskip}
\blockanchor{p108-b08}
\majorquestion{一}(15 分) 判断极限 $\lim_{n\to \infty}\sum_{k = 1}^{n}\ln \left(1 + \frac{k}{n^2}\right)$ 是否存在？若存在，求出极限值.\par
% p108-b09; source PDF page 108
\par\Needspace{5\baselineskip}
\blockanchor{p108-b09}
\majorquestion{二}(15 分) 设 $f(x)$ 在 $[0,1]$ 上连续, 在 $(0,1)$ 上可导, $f(0) = 0, f(1) = 1$, 证明: 在 $(0,1)$ 内存在不同的 $\xi, \eta$, 使得 $f'(\xi)f'(\eta) = 1$.\par
% p108-b10; source PDF page 108
\par\Needspace{5\baselineskip}
\blockanchor{p108-b10}
\majorquestion{三}(15 分) 设 $f(x)$ 在 $[1, +\infty)$ 上连续, 且 $\lim_{x \to +\infty} [f(x) - 3\sqrt{x}] = 7$. 判断函数 $f(x)$ 在 $[1, +\infty)$ 上是否一致连续, 并说明你的理由.\par
% p108-b11; source PDF page 108
\par\Needspace{5\baselineskip}
\blockanchor{p108-b11}
\majorquestion{四}(15 分) 设 $f(x)$ 在 $[a, b]$ 上有定义, 并且在 $[a, b]$ 上每一点的极限都存在且为 0, 证明: $f(x)$ 在 $[a, b]$ 上可积, 并求 $\int_{a}^{b} f(x) \mathrm{d}x$.\par
% p108-b12; source PDF page 108
\par\Needspace{5\baselineskip}
\par\Needspace{6\baselineskip}
\blockanchor{p108-b12}
\majorquestion{五}(15 分) 设 $f(x)$ 在 $[a, b]$ 上连续, 且 $\min_{x \in [a, b]} f(x) = 3$, 计算\par
\nopagebreak[4]
% p108-b13; source PDF page 108
\blockanchor{p108-b13}
\displayformula{\lim _ {n \rightarrow \infty} \left[ \int_ {a} ^ {b} \frac {1}{(f (x)) ^ {n}} \mathrm{d} x \right] ^ {\frac {1}{n}}.}
% p108-b14; source PDF page 108
\par\Needspace{5\baselineskip}
\par\Needspace{6\baselineskip}
\blockanchor{p108-b14}
\majorquestion{六}(15 分) 设 $a_{n} = \int_{0}^{\frac{\pi}{4}} \tan^{n} x \, \mathrm{d}x$, 计算 $\sum_{n=1}^{\infty} \frac{a_{n} + a_{n+2}}{n}$ 的值, 判断级数 $\sum_{n=1}^{\infty} \frac{a_{n}}{\sqrt{n}}$ 是否收敛? 并说明理由.
\majorquestion{七}(15 分) 设 $F(x, y)$ 在 $(x_{0}, y_{0})$ 某邻域内有二阶连续的偏导数, 且\par
\nopagebreak[4]
% p108-b15; source PDF page 108
\blockanchor{p108-b15}
\displayformula{F (x _ {0}, y _ {0}) = 0, F _ {x} ^ {\prime} (x _ {0}, y _ {0}) = 0, F _ {x x} ^ {\prime \prime} (x _ {0}, y _ {0}) F _ {y} ^ {\prime} (x _ {0}, y _ {0}) \neq 0.}
% p108-b16; source PDF page 108
\blockanchor{p108-b16}
证明: 由方程 $F(x, y) = 0$ 在 $(x_0, y_0)$ 某邻域内确定的隐函数 $y = y(x)$ 在点 $x_0$ 处取极值.
\majorquestion{八}(15 分) 求平面 $x + 2y = 1$ 与曲面 $x^2 + 2y^2 + z^2 = 1$ 的交线上距离原点最近的点的坐标.\par
% p108-b17; source PDF page 108
\par\Needspace{5\baselineskip}
\blockanchor{p108-b17}
\majorquestion{九}(15 分) 求曲面 $(x^{2} + y^{2} + z^{2})^{3} = a^{3}xyz (a > 0)$ 所围立体的体积.\par
% p108-b18; source PDF page 108
\par\Needspace{5\baselineskip}
\blockanchor{p108-b18}
\majorquestion{十}(15 分) 证明: $F(\lambda) = \int_{0}^{+\infty} \frac{1 - e^{-\lambda t}}{t} \cos t \, \mathrm{d}t$ 在 $[0, +\infty)$ 上连续, 在 $(0, +\infty)$ 内可导.\par
